\documentclass[11pt]{article}
\usepackage[margin=1.2in]{geometry}
\usepackage{amsmath,amssymb,amsthm}
\usepackage{hyperref}
\hypersetup{colorlinks=true, linkcolor=blue, urlcolor=blue}
\usepackage{parskip}
\emergencystretch=3em

\newcommand{\R}{\mathbb{R}}
\newcommand{\code}[1]{\texttt{#1}}

\title{Tide retrospective: the forward theorems\\
\large the germbij forward direction closes}
\author{Claude (Anthropic), with GPT-5.6 Sol consults}
\date{2026-08-10}

\begin{document}
\maketitle

\begin{abstract}
The public face of the forward-expansion programme, and its close:
on a \code{ForwardExpansionDomain} --- a nondegenerate local minimum
with $C^{N+2}$ regularity and the Peano remainder --- the rescaled
posterior moment of every continuous observable of polynomial growth
is an order-$N$ asymptotic polynomial at $0^+$
(\code{rescaledMoment\_hasExpansion}). The coefficients are the
division of the numerator expansion by the partition expansion, whose
constant coefficient $\int e^{-T_2}$ is positive; the monomial
instances are exactly the moment families that the merged
inverse-half theorems consume; and moment families agreeing to
$o(q^N)$ have equal coefficients through order $N$. Together with the
inverse half --- superpolynomially-close localized moment families
determine the location, the full positive-order jet, and for analytic
losses the germ --- this completes the germbij note's Theorem~3.1 in
the nondegenerate case: expansion data and Taylor data determine each
other, in both directions, in Lean.
\end{abstract}

\section{Setting}

Stage 7 of the seven-stage design consult, which ran start to finish
as a single stacked arc: the expansion predicate and coefficient
uniqueness; the mesoscopic window and its tail calculus; the
one-field Peano mixin and the exact exponent split; the graded
exponential; the coefficient functions; the scalar bounds; the
quadratic bridge and Gaussian absorption; the $q$-uniform majorant;
the window dominated convergence; the outer tails and the numerator
packaging; asymptotic division; and now the public statements. Each
stage's interface was consumed by the next with essentially no
rework --- the two architecture consults (the programme design and
the stage-5 shape) did what they were commissioned to do.

\section{The thing formalised}

\begin{sloppypar}
\code{numeratorCoeff\_one\_zero\_pos}: the partition constant
coefficient is $\int e^{-T_2} > 0$, from the zeroth coefficient
function being constantly one, the support of a positive integrand
being everything, and the volume measure being nonzero.
\code{momentCoeff} and \code{rescaledMoment\_hasExpansion}: the
division of the two numerator expansions, applied at the seabed's
\code{rescaledMoment}. \code{rescaledMoment\_expansion\_exists} (the
existence form), \code{monomial\_moment\_hasExpansion} (the
\code{monomialTest} instances), and
\code{momentCoeff\_eq\_of\_isLittleO} (the comparison face, through
stage A1's uniqueness).
\end{sloppypar}

\section{Where progress stalled}

\begin{sloppypar}
One rename: \code{IsLittleO.neg} is spelled
\code{IsLittleO.neg\_left}. The tide compiled on the first fix
round --- fitting for a stage that was designed to be pure
composition.
\end{sloppypar}

\section{Mathlib lookup versus new mathematics}

\begin{sloppypar}
Lookup: the support characterization of positive integrals
(\code{integral\_pos\_iff\_support\_of\_nonneg}),
\code{Set.eq\_univ\_of\_forall}, \code{Function.mem\_support},
\code{Measure.measure\_univ\_pos}, \code{NeZero.ne}. New: none. The
mathematics of this tide is the programme it closes.
\end{sloppypar}

\section{What the next tide inherits}

Process work first: the blue-dot \code{\textbackslash leanref} marker
pass on the germbij note for the forward theorems, the SRI docs sync,
and the commit-pin bump. Mathematically, what remains of the germbij
note beyond the user's stated core: the semi-quasi-homogeneous
constructive class, and --- if explicit coefficient values are ever
wanted rather than existence --- the connection of
\code{momentCoeff} to Wick/Isserlis data through the Gaussian moment
integrals already in the seabed.
\end{document}
